\section{Persistent marginal information}\label{sec:persistence}

All logarithms are base two. For a finite random variable $W$,
$H_\infty(W)=-\log_2\max_w\Pr(W=w)$. A worst-case conditional bound
ranges over attained conditioning values of positive probability. It is
stronger than a bound on the average guessing probability. Neither notion
implies independence of two work words.

\begin{lemma}[One subset for every context]\label{lem:context-pruning}
\formal{ContextPruning.lean}{context_subset,context_half_mass,context_conditional_bound}
Let $(\Omega,\mu)$ be a finite probability space, $\mu(B)=p>0$, and let $Z_1,\ldots,Z_\ell$
be context variables. For $\theta>0$ there is $D\subseteq B$ with
$\mu(D)\ge p-\ell\theta$ and
\[
 \mu(D\cap\{Z_i=z\})\ge\theta\mu(Z_i=z)
\]
for every context attained on $D$. If, for every $i,v,w,z$,
$\mu(B,W_{v,i}=w,Z_i=z)\le c\mu(Z_i=z)$, with $c\ge0$ and $\ell\ge1$, choosing
$\theta=p/(2\ell)$ gives on this one $D$ of mass at least $p/2$
\[
 \Pr_D(W_{v,i}=w\mid Z_i=z)\le 2c\ell/p.
\]
\end{lemma}
\begin{outline}
For $E\subseteq B$, let $S_i(E)$ be the union in $\Omega$ of all $i$-context
fibres meeting $E$. Maximize the finite potential
$\Phi(E)=\mu(E)-\theta\sum_i\mu(S_i(E))$. For a maximizer $D$,
$\mu(D)\ge\Phi(D)\ge\Phi(B)\ge p-\ell\theta$. If an attained fibre
$F=D\cap\{Z_i=z\}$ violated the asserted lower bound, removing $F$ would
lose $\mu(F)$ but remove the whole penalty $\theta\mu(Z_i=z)$ for this
context. All other penalties can only decrease, contradicting maximality.
Finally divide the joint envelope by this fibre lower bound. The global
conditioning factor $\mu(D)$ cancels. There is no extra factor for the
number of word variables sharing each context partition.
\end{outline}

\begin{theorem}[Persistent information beyond the strict output prefix]\label{thm:conditional-information}
\partformal{PersistenceSource.lean}{ConditionalPersistence.of_small_thin_source}{the native catalogue counting theorem. Given its exact thin-source cardinality bound, including its factor two, Lean constructs the original uniform law, escaping set, common live interval, retained set, conditional whole-word bound and uniform threshold.}
Fix $k\ge2,a,C$ and the constants $\rho,\eta$ of \cref{thm:catalogue}.
There are constants $\alpha,\beta,p_0>0$ such that, for every sufficiently
large $d$, every $b\ge1$ and machine with $q\le\rho P/2$, the following
holds. If $S$ is a timely correct set of size at least
$2^{(1-\eta/2)P}$, there are two fixed tapes, one fixed interval of at least
$\beta P/4$ live configuration times, and one $D\subseteq S$ with
$|D|/|S|\ge p_0/2$, such that at every time of the interval both whole
work words have worst-case min-entropy at least $\alpha P/2$ conditional
on the strict output prefix, under the uniform law on $D$. On every selected
input, at least $\alpha P$ output bits remain undelivered throughout.
\end{theorem}
\begin{informalproof}
Set $h_0=\rho/[4(k+1)]$, $A_0=2^a$,
$\alpha=h_0/4$, $\beta=h_0/[4\log_2(6A_0)]$,
$R=\lfloor\beta P\rfloor$, $\ell=\lfloor R/2\rfloor$ and
$K_0=\lceil2^{\alpha P}\rceil$. Fix the original uniform law $\mu$ on $S$.
A one-tape vertex is its entire word on $[-T_*,T_*]$ and its head position.
Writing one symbol and making one unit move gives at most $3A_0$ outgoing
and $3A_0$ incoming neighbours. Hence an undirected radius-$R$ ball has
at most $(6A_0)^{R+1}$ vertices.

For each time $t$, prefix $z$ and tape, take its $K_0$ most probable positive
conditional vertices, or all if fewer exist. Outside this set each conditional
atom has mass at most $1/(K_0+1)$. Define a public catalogue at $(t,y)$ by thickening the top sets at
$u=(t-l)_+$ and prefixes $y[0:(|y|-r)_+)$ for $0\le l,r\le R$.
Use radius-$R$ balls and map each vertex to its whole word, blank outside
$[-T_*,T_*]$. This square overcount includes every actual historical
prefix: at lag $l$, the prefix loses at most $l$ bits. Forgetting the head
cannot increase the catalogue size, so its size is at most
\[
 K=(R+1)^2K_0(6A_0)^{R+1}.
\]
Every actual historical prefix occurs among these candidates because at most
one bit is appended per transition. Its logarithm, including the ceiling,
is at most $(h_0/2)P+O(\log P)\le h_0P$ eventually. The catalogues are
fixed from the original law before applying \cref{thm:catalogue}.

That theorem implies that more than half of $S$ have a reached time $t$ and
two words escaping their current catalogues. Every vertex of both tapes in
the preceding $R$ steps then avoids its corresponding conditional top set.
The witness time exceeds $R$: otherwise the certain initial blank vertex
and its actual path put every word in its catalogue. Choose
$I_j=\{j\ell,\ldots,(j+1)\ell-1\}$ with $j=\lfloor t/\ell\rfloor-1$.
It lies in $[t-R,t)$ and is strictly before halt, even if $t$ itself is the
halting configuration. There are at most
$M_0=\binom k2(\lceil8C/\beta\rceil+1)$ pair/grid choices. One cell $B$
of this partition has mass $p\ge p_0=1/(2M_0)$.

At each time in this fixed interval, both tapes avoid their conditional top
sets throughout $B$. Forgetting the head position costs at most
$n=2T_*+1$ atoms, giving the original-law envelope
\[
 \mu(B,W_{v,t}=w,Z_t=z)\le \frac{n}{K_0+1}\mu(Z_t=z).
\]
Apply \cref{lem:context-pruning} to the $\ell$ time-prefix partitions. One
$D$ of mass at least $p_0/2$ has every conditional word atom at most
$2n\ell/[(K_0+1)p_0]$. Its negative logarithm is
$\alpha P-O(\log P)\ge\alpha P/2$. A positive attained vertex outside
the top set also witnesses at least $K_0+1$ inputs in its original
time/prefix fibre. Correctness fixes $|z|$ distinct raw bits, so that actual
fibre has size at most $2^{P-|z|}$. Thus $P-|z|\ge\log_2(K_0+1)\ge\alpha P$.
Adaptive output timing can further shrink the fibre; equality with
$2^{P-|z|}$ is not assumed. All bounds are uniform for $P\ge d^2$.

The formal endpoint retains the native catalogue bound
$2^{1+(1-\eta)P}$ as its sole unproved input. The source-size bound and
$\eta P/2\ge2$ leave at least half the original uniform source outside
that thin set. The shared
\lean{LeanMachine.FiniteSourceEscape.half_mass_pair_escape} constructs
that event and its actual two-tape witnesses. The remaining selection,
probability and rate estimates are formalized, with the dimension threshold
preceding the width, machine and source quantifiers.
\lean{ConditionalPersistence.of_half_mass_escapes} also works under any
strictly positive finite source law satisfying the literal escape condition.
The context-pruning step uses one retained set for every selected time and
tape. The prefix-fibre bound follows from actual native correctness; no
full-size fibre, independent tape words or entropy lower bound is assumed.
\end{informalproof}

Taking no prefix condition in the same argument gives the ordinary
persistent-marginal theorem. The stronger conditional result still permits
the two words to be identical. It cannot be summed over times as repeated
fresh information.

\subsection{Future successful reads and remote returns}

Define the future-readable word at source time $t$ by retaining the current
symbol exactly at cells visited in successful steps $u\ge t$ before halt;
at and after halt it is blank. Call it $L_h(t)$. The unused halt-test tuple
does not count as a future read. Its nonempty future address set is an integer
interval, by unit motion.

\begin{theorem}[Source future-readable catalogue bound]\label{thm:live-catalogue}
\notformalized
For fixed $k,a,C$ there are $h,e,\rho'>0$ such that, uniformly in $b$ at
large $d$, every public time-indexed family of catalogues of size at most
$K$ with $\lceil\log_2K\rceil\le hP$ has fewer than $2^{(1-e)P}$ timely
correct inputs for which at every source time at most one $L_h(t)$ lies
outside its catalogue. One may allow $q\le\rho'P$. No online access to
future information is given to the source.
\end{theorem}
\begin{informalproof}
Use the aligned bit simulator of \cite{HighRateStorage2026} with block size
$s\ge1$ and physical batch length $B=10s+1$. Its literal tape word at
time $Bj$ equals the source word at time $js$. Each batch reads/writes a
public $3s$-cell window around its anchor; outside that window the word remains
the source boundary word throughout the batch. Its positions at successful
time $u$ differ by at most $2s$ from the source at
$\min(s\lfloor u/B\rfloor,t_X-1)$, including the final batch and stationary
drain. These identities follow from the simulator's boundary invariant and
its actual sweep/carry rule, with source offset in $[0,s-1]$ and physical
offset in $[-s,2s-1]$. The possible terminal replacement $t_X\mapsto t_X-1$
costs one unit.

Fix one full simulator trace of successful horizon $U$. Hardwire its last
successful visit flags and clock in a native eraser: simulate each old tuple,
but replace a work write by blank at the last visit to that cell. Force halt
at $U$, because the erased final tuple need not itself trigger the source
halt rule. On this trace fibre, induction shows its word is exactly the
simulator word masked to future successful visits, with unchanged output and
transition count. It has at most $1+Q_{\mathrm{sim}}(U+1)$ states.

For an ordinary batch $j<\lfloor t_X/s\rfloor$, let $J$ be this trace's
public future simulator interval and $I_X$ the source future interval at $js$.
The shadow gives $J\subseteq I_X+[-2s,2s]$. Remove the first and last $2s$
cells of $J$. The remainder is in $I_X$. Outside these at most $4s$ cells
and the current $3s$ write window, the eraser word is exactly the source
future-readable word masked to $J$. Thus a source catalogue induces one of
size at most $KA^{7s}$. All heads use the same source time $js$. In the final
batch and drain the entire future word is supported in the public $3s$
window; at halt it is blank. This covers divisible source runtimes and $s=1$.
No source future-interval endpoints are supplied as free metadata.

Set $C_N=11C+1$ and first choose the catalogue constants of
\cref{thm:catalogue} for $(k,a,C_N)$, independently of $s$. Write them
$\rho_N,\eta_N$ and $h_N=\rho_N/[4(k+1)]$. The published successful-input
simulator bound gives at most $2^{a_0P/s}$ full traces, where
$a_0=(2C+1)\log_2(1+2\cdot3^{k+1})$, for large $P$. Choose fixed $s$ with
$a_0/s\le\eta_N/4$. Its control cardinality is at most $QD_s$, with
\[
 D_s=(11s+2)s^{k+2}3^{3s+k+1}(2^a)^{3sk}2^{2s+2}.
\]
This counts the actual phase, offsets, raw/work windows, carry directions,
pending/scheduled bits and stopped flag. The eraser consequently costs only
$q+\lceil\log_2D_s\rceil+\lceil\log_2(U+1)\rceil+1$ bits. Take
$\rho'=\rho_N/4$ and $h=h_N/2$; fixed and logarithmic overhead eventually
fit the remaining reserves. Apply the ordinary time-indexed catalogue bound
separately on every trace fibre and sum, retaining its factor two. The total
is less than $2^{1+(1-\eta_N+a_0/s)P}<2^{(1-\eta_N/2)P}$ eventually.
Thus $e=\eta_N/2$ works. No largest-fibre selection loses the original set.
\end{informalproof}

\begin{corollary}[Persistent information on original remote returns]\label{cor:remote-information}
\partformal{PersistenceFutureSource.lean}{PersistenceFutureSource.of_thin_source_bound}{the source future-readable catalogue counting theorem. From its explicit thin-source bound, Lean derives the original uniform law, one common event and live interval, whole-word and remote-word atom bounds, actual successful returns and a threshold uniform over positive widths.}
For fixed $k\ge2$ and resource-dependent $\alpha,\beta,p_0>0$, every sufficiently large timely
correct set $S$ of size at least $2^{(1-e/2)P}$ contains a subset of relative
size at least $p_0$, two fixed tapes and a common interval of at least
$\beta P$ live times on which both future-readable words have min-entropy
at least $\alpha P$. The same holds, with smaller constants, after blanking
a fixed radius-$\alpha'P$ window about each current head. Every nonblank
cell of the latter word lies on an actual macroscopic return straddling
that checkpoint. These are marginal, not prefix-conditioned, claims.
\end{corollary}
\begin{informalproof}
The pair $(L_h(t),y_h(t))$ follows the same bounded-degree one-tape graph used
above: only the old scanned cell changes, retaining its write if revisited
later and otherwise becoming blank. Thicken the $K_0$ most likely vertices
at all nearby times under the original uniform law on $S$, now without prefix
indices. At radius $R=\Theta(P)$ the catalogue size is
$(2R+1)K_0(6A_0)^{R+1}$, which fits \cref{thm:live-catalogue} by choosing
small exponential rate and radius. Escape on more than half of $S$ gives a
fixed pair/grid interval of constant relative mass, exactly as in the preceding
proof. Full vertex atoms there are at most $1/(K_0p_0)$; forgetting the head
costs at most $\log_2(2T_*+1)$ bits.

Masking a radius-$H$ window has fibres of size at most
$(2T_*+1)A_0^{2H+1}$ on full vertices. Choose a sufficiently small fixed
$H/P>0$; linear min-entropy survives. A nonblank retained cell was written
at some previous visit and has a future successful visit. Its last past and
next future visits are consecutive and its distance from the current head
exceeds $H$, so the checkpoint witnesses the claimed return. All selected
times are before halt. Neither this graph argument nor the eraser gives the
original source a future-read oracle or proves that its later read matters
semantically to the output.

The formal proof uses one finite observation window $[-T_*,T_*]$ for the
whole source law, while retaining each input's own first halt time in its
future-readable word. Frozen steps after that halt are never future reads.
\lean{PersistenceFutureRates.of_half_mass_escapes_all_widths} proves the
selection and quantitative bounds under arbitrary finite laws; the stated
uniform-source result derives its half-mass escape event from the thin-source
count. The native eraser and source-catalogue theorem remain the only
unformalized dependency of this consequence.
\end{informalproof}
